% Generated draft: requires mathematical and source review.
\documentclass[11pt]{article}
\usepackage[margin=25mm]{geometry}
\usepackage{fontspec}
\setmainfont{Latin Modern Roman}
\usepackage{amsmath,amssymb,mathtools,unicode-math}
\setmathfont{Latin Modern Math}
\usepackage{enumitem,needspace,xurl,hyperref}
\setlength{\parindent}{0pt}
\setlength{\parskip}{5pt}
\setlength{\emergencystretch}{3em}
\newcommand{\N}{\mathbb{N}}
\newcommand{\Z}{\mathbb{Z}}
\newcommand{\Q}{\mathbb{Q}}
\newcommand{\R}{\mathbb{R}}
\newcommand{\C}{\mathbb{C}}
\newcommand{\F}{\mathbb{F}}
\newcommand{\A}{\mathbb{A}}
\newcommand{\PP}{\mathbb{P}}
\newcommand{\E}{\mathbb{E}}
\newcommand{\eps}{\varepsilon}
\newcommand{\dd}{\,\mathrm{d}}
\begin{document}
\section*{UnsolvedMath ID: 30004855}
\textbf{Source problem code:} catalogue entry
\subsection{Definitions and mathematical statement}
Let \(n\geq 1\) and let \(f\in\Z[x_1,\ldots,x_n]\) have total degree \(d>1\). Write \(f_d\) for the sum of the monomials of \(f\) of total degree \(d\). The source denotes by \(s\) the dimension of the critical locus of \(f_d\). For a specified field \(K\), the corresponding affine critical locus is
\[
\operatorname{Crit}_K(f_d)=\left\{a\in\overline K^{\,n}:\frac{\partial f_d}{\partial x_i}(a)=0\text{ for every }1\leq i\leq n\right\}.
\]
The supplied record does not specify the ground field for this dimension. It also does not define \(E_f(p,m)\), including its normalization, summation domain, or exponential character. Accordingly, \(s\) and \(E_f(p,m)\) below retain their intended source meanings; the supplied evidence does not determine a fully specified exponential-sum conjecture.

\textbf{Extracted conjectural bound, subject to these unresolved definitions.} For every such fixed \(f\), every real \(\varepsilon>0\), and every prime \(p\), there exists a real constant \(C_{f,p,\varepsilon}>0\), independent of \(m\), such that
\[
\lvert E_f(p,m)\rvert\leq C_{f,p,\varepsilon}\,p^{m\left(-\frac{n-s}{d}+\varepsilon\right)}
\qquad\text{for every integer }m\geq 1.
\]
Moreover, for each fixed \(f\) and \(\varepsilon>0\), there exists a threshold \(P_{f,\varepsilon}\) such that for every prime \(p>P_{f,\varepsilon}\) and every integer \(m\geq1\),
\[
\lvert E_f(p,m)\rvert\leq p^{m\left(-\frac{n-s}{d}+\varepsilon\right)}.
\]
The subscript \(f\) makes explicit the permitted dependence on the fixed polynomial; the source writes \(C_{p,\varepsilon}\). No uniformity across different polynomials is stated.

\textbf{Supplied qualifications.} The report passage describes this as a finite-ring analogue of Deligne's theorem on exponential sums over finite fields and presents its exponent as \(\sigma<(n-s)/d\). Its unidentified speaker reports joint work with Cluckers establishing the displayed conjecture for polynomials in at most four variables or satisfying an unspecified non-degeneracy condition. The preceding discussion reports joint work with Veys on the rational-singularity case in at most three variables, using a classification of rational singularities; the antecedent of ``this conjecture'' in that preceding discussion is not fully supplied. These remarks do not impose a rational-singularity hypothesis on the displayed general conjecture.
\subsection{Short English statement}
For a fixed integer polynomial, the proposed bound controls the source's finite-ring exponential sums at every prime-power level, with constant \(1\) for sufficiently large primes. The sum itself and the field used to define the critical-locus dimension are missing from the record, so this remains a statement-only draft requiring mathematical review.
\subsection{Sources}
\begin{itemize}
\item Local registry: \texttt{data/unsolvedmath/problems.json}, permanent UnsolvedMath ID \(30004855\), exact source code \texttt{catalogue entry}, title ``Igusa Bounds for Finite-Ring Exponential Sums.'' The supplied record attributes the extraction to the Oberwolfach Reports Open Problems collection; no proposer is named.
\item Supplied report citation: \emph{Singularities (2022)}, Oberwolfach Reports, DOI \texttt{10.4171/owr/2021/47}, pp.\ 2573--2639. \url{https://doi.org/10.4171/owr/2021/47}. The record lists proposed year \(2021\). The passage mentions joint work with Veys and with Cluckers but does not identify its speaker.
\item Supplied literature source: Cluckers--Mustaţă--Nguyen, ``optimal nonrational estimates'' (the supplied descriptive label). \url{https://doi.org/10.1017/fmp.2019.3}.
\item Supplied literature source: ``OWR report recording the remaining cases.'' \url{https://doi.org/10.4171/owr/2021/47}.
\item Literature assessment dated exactly \texttt{2026-08-22}: the supplied curation labels the problem partially solved and reports Igusa-type bounds for nonrational singularities, a rational-singularity case in at most three variables, and the stated exponent in at most four variables or for nondegenerate polynomials, with the general conjecture unresolved. This is supplied dated curation, not an independently verified literature assessment. The record itself requests independent verification of time-sensitive status and cited claims.
\end{itemize}
\end{document}
